\section{Limitations and Future Work}
\label{sec:limitations}

Our results should be interpreted within the bounded settings in which the
system and studies were evaluated. The formative study examined
\DataNeeded{CHARACTERIZE THE PARTICIPANT POPULATION AND RANGE OF REPAIR
EPISODES}, and the controlled study placed participants in the role of a
workflow owner during two short, consequential but simulated tasks. This
design allowed us to compare conditions under matched evidence, but it does
not reproduce the accumulated context, organizational pressure, or delayed
consequences of maintaining an agent over weeks or months. Participants'
authority in the study was assigned by the task; in practice, the person
editing or approving a skill may not be the person most affected by its
behavior. Longitudinal deployments should examine how behavioral scopes and
release records evolve through repeated breakdowns and handoffs among owners,
maintainers, and affected stakeholders.

The controlled study used two fixed domains, one model configuration, typed
tool interfaces, and sandboxed world states. These choices provided
reconstructable outcomes and prevented study actions from affecting real
accounts. They also limit ecological and domain generality. Travel recovery
and expense review contain explicit decisions, constraints, and side effects;
skills for creative, interpersonal, or open-ended analytical work may not
admit equally clear fact or trace criteria. Future work should study domains in
which behavior must be reviewed through human judgment, contested values, or
longer temporal horizons, and should test whether the interaction remains
useful across models and changing tool environments.

Our scenario packs supplied a bounded set of realistic contrastive cases. In
an open-world product, relevant cases may be missing from logs, expensive to
reconstruct, or impossible to execute safely. Model-generated cases may also
reflect the model's own assumptions and narrow the policy space presented to
the owner. \systemname{} therefore treats suggested cases as proposals and
retains their provenance, but the present study cannot establish the coverage
of this process. Retrieval from historical incidents, collaborative case
authoring, and methods for identifying absent stakeholder perspectives are
important complements to automated contrast generation.

The evidence mechanisms have corresponding technical limits. Repeated
executions characterize behavior only under the selected cases, artifact
versions, model settings, and tool snapshots. A fixed execution budget may
miss low-frequency failures, while model or service updates can invalidate an
earlier comparison. Deterministic evaluators improve reconstructability but
cover only behavior represented in the trace or world state. Targeted
instruction interventions provide conditional sensitivity evidence; they do
not identify unique causality, exhaust interactions among instructions, or
guarantee that a candidate will generalize. Behavioral patch review should
therefore be understood as bounded review evidence, not verification or safety
certification.

The user study compares an integrated workspace with a capability-matched chat
interface. This design tests the behavioral patch review model as a whole but
cannot attribute an observed benefit to persistent scope, cross-highlighting,
the structured diff, or side-by-side execution evidence individually. Factorial
or ablation studies could examine which representations are necessary under
different repair complexity. The within-participant design also introduces
possible learning and carryover despite counterbalancing, and the two domains
are fixed task contexts rather than samples from a population. We report order
and domain effects but avoid claiming domain-independent effect sizes.

Several outcomes measure different notions of repair quality. Scope adherence
captures whether a candidate respects the participant's prospective
commitments, not whether those commitments constitute an objectively correct
policy. Hidden cases and preregistered criteria provide complementary evidence
about transfer and consequential failures, but they necessarily encode
researcher judgments about the scenarios. Likewise, decision quality evaluates
consistency with visible mismatches and evidence gaps; it does not imply that
release is mandatory whenever those blockers are absent. Future studies should
examine collective review, disagreement among legitimate authorities, and
decisions for which several policies are defensible.

Finally, making a release auditable does not make it legitimate. Execution
traces and historical cases can contain sensitive personal or organizational
information, and preserving provenance can conflict with data minimization or
deletion requirements. A workflow owner may also lack authority to define
behavior for people affected by the agent. Production systems will require
access controls, retention policies, appeal mechanisms, and ways to represent
stakeholders who are not present during repair. Behavioral patch review offers
a structure for exposing commitments and evidence, but broader governance must
determine who may author those commitments and who bears the consequences of a
released change.
